% Appendix. Fragment, \input by spacenav.tex inside \appendix.
\section{Observation channels}
\label{app:channels}

Section~\ref{sec:instruments} summarizes the observation-channel result;
this appendix gives the channels, the redundancy argument and the width
experiment in full. The four channels added to the observation are $\eta$; a log-compressed authority $A$; a log-compressed
demand $\dd$; and the local tidal stretching rate (the square root of the
tidal gradient, i.e. how fast nearby trajectories are pulled apart). The
channels can be zeroed individually or together with the observation shape
unchanged, so that an ablated run stays comparable with a full one, and
toggling them was verified not to perturb the dynamics.

Supplying all four channels costs the agent between 3.1 and 10.1 percentage
points of final success depending on the tier (Table~\ref{tab:instr};
Figure~\ref{fig:ablation}(a); Welch's $t$-test on seed means, $n = 5$ per
arm, $p = 0.021$, $0.0088$, $0.029$ and $0.0017$ on the easy, medium, hard
and blind suites, with effect sizes between $1.8$ and $3.0$). The quoted
uncertainties are training-seed standard deviations, which are the relevant
ones; the episode-level binomial intervals are far narrower and would
overstate the confidence.

The observation already carries the local gravity vector (hence
$\lVert g \rVert$), the product of fuel and engine acceleration (hence
$\Delta v_{\mathrm{budget}}$), and the fraction of the episode remaining
(hence $T$), so that $A = \Delta v_{\mathrm{budget}} / (\lVert g \rVert T)$ is
an algebraic combination of three quantities the agent already has, partly
duplicating an existing ratio of remaining to required $\Delta v$, and
$\dd = \eta / A$ is redundant given $\eta$ and $A$; by contrast $\eta$
requires the dominant-body decomposition, which is not otherwise computable
from the observation, and the stretching rate is not present in any form
(it is also a good proxy for the quantity one would actually want, since on
the four levels of Figure~\ref{fig:phase} its log–log correlation with the
measured 20\,s Lyapunov exponent is $0.96$--$0.98$).
Splitting the panel along that line puts the entire cost on the redundant
half (Table~\ref{tab:instr}; Figure~\ref{fig:ablation}(b)). Supplying only
$\eta$ and the stretching rate is statistically indistinguishable from
supplying nothing ($p = 0.40$ to $0.90$ across suites), supplying only $A$
and $\dd$ is indistinguishable from supplying all four ($p = 0.59$ to
$0.81$), and the contrast between the two halves is significant on every
suite ($p \le 0.0064$, effect sizes $2.4$ to $3.6$).



In this environment, at this network size and step budget, appending
features that are algebraic combinations of what the observation already
contains measurably degrades a policy, while appending features that carry
physics not otherwise available does not, and there is no evidence that the
informative channels help ($\eta$ and the stretching rate are nominally
above the no-channel arm on three suites and below it on one, none
significantly). Two accounts of the mechanism predict different behavior
under scale: four extra input dimensions against 285,322 parameters is a
capacity account, under which a wider network should absorb the redundant
channels, while a ratio with a division in it is heavy-tailed even after log
compression, which is a conditioning account under which the cost should
persist at any width. Doubling the embedding width to 256 (698,634
parameters, three seeds per arm, same budget) decides between them: the
cost falls to $-0.3$, $+2.4$, $+0.8$ and $+2.7$ points on the four tiers
(Welch's $t$, $p = 0.74$, $0.18$, $0.38$ and $0.026$), the all-four arm at
the wider network ($90.1 \pm 0.9$, $41.2 \pm 0.4$, $11.4 \pm 1.2$,
$29.6 \pm 0.5$) matches the $\eta$,stretch arm at either width, and the
$\eta$,stretch arm itself gains little from the extra width ($-0.6$, $+1.7$,
$+1.1$ and $+3.1$). The cost was capacity. Thus, an observation channel's
value to a learner depends on whether it is new rather than on whether it is
true, and at a network size where capacity is scarce a channel that does not
change the policy costs capacity without buying anything, which is
consistent with the view of a policy as a compression of its inputs
\cite{lai2024,kaelbling1998}; at a size where capacity is not scarce the
redundant channels are merely ignored, and the wider network spends its
extra parameters on tolerating them rather than on the task.
